\section*{अध्याय \ref{ch:g11:seismic-structure} --- भूकंपीय तरंगें पृथ्वी का भीतरी भाग प्रकट करती हैं}

\begin{solution}{exo:g11:seismic-structure:1}
\omterm{def:g11:seismic-structure:body-waves}{P तरंग} चट्टान को संपीडित करती और खींचती है, और चट्टान गति की दिशा के अनुदिश
चलती है; \omterm{def:g11:seismic-structure:body-waves}{S तरंग} उसका अपरूपण करती है, और चट्टान उस दिशा के समकोण पर
चलती है। \omterm{def:g11:seismic-structure:body-waves}{P तरंगें} ठोसों और द्रवों को पार करती हैं; \omterm{def:g11:seismic-structure:body-waves}{S तरंगें} केवल ठोसों को। \omterm{def:g11:seismic-structure:body-waves}{P
तरंग} पहले दर्ज होती है, क्योंकि दोनों में वही अधिक तेज़ है।
\end{solution}

\begin{solution}{exo:g11:seismic-structure:2}
$40 \times 111.2 = \qty{4448}{km}$, लगभग \qty{4450}{km}। $8900/111.2
\approx \ang{80}$।
\end{solution}

\begin{solution}{exo:g11:seismic-structure:3}
\qty{456}{s} का अर्थ 7~मिनट 36~सेकंड है; \qty{822}{s} का अर्थ 13~मिनट
42~सेकंड। S--P देरी $822 - 456 = \qty{366}{s}$ है, यानी 6~मिनट
6~सेकंड।
\end{solution}

\begin{solution}{exo:g11:seismic-structure:4}
$\sin i_c = 6.0/8.0 = 0.75$, इसलिए $i_c \approx \ang{48.6}$।
\end{solution}

\begin{solution}{exo:g11:seismic-structure:5}
बाह्य क्रोड $5150 - 2891 = \qty{2259}{km}$ मोटा है। क्रोड की त्रिज्या
$6371 - 2891 = \qty{3480}{km}$ है, और $3480/6371 \approx 0.55$: पृथ्वी की त्रिज्या के
आधे से थोड़ा अधिक।
\end{solution}

\begin{solution}{exo:g11:seismic-structure:6}
6~मिनट 6~सेकंड का अर्थ \qty{366}{s} है: भूकंप \ang{40} दूर है, यानी
$40 \times 111.2 \approx \qty{4450}{km}$। एक अकेली देरी केवल
दूरी देती है: अधिकेंद्र स्टेशन के चारों ओर उस त्रिज्या के वृत्त पर कहीं भी हो
सकता है। तीनों से सही दूरी वाला एकमात्र बिंदु ढूँढने के लिए
दो और स्टेशन चाहिए।
\end{solution}

\begin{solution}{exo:g11:seismic-structure:7}
$i_c \approx \ang{48.6}$ और $\cos i_c \approx 0.661$। सीधी तरंग: $200/6.0
\approx \qty{33.3}{s}$। \omterm{ex:g11:seismic-structure:sp}{शीर्ष तरंग}: $200/8.0 + 2 \times 35 \times 0.661/6.0
\approx 25.0 + 7.7 = \qty{32.7}{s}$। \omterm{ex:g11:seismic-structure:sp}{शीर्ष तरंग} पहले पहुँचती है, लगभग
\qty{0.6}{s} पहले।
\end{solution}

\begin{solution}{exo:g11:seismic-structure:8}
$\tan i_c \approx 1.134$: क्रांतिक दूरी $2 \times 35 \times 1.134
\approx \qty{79}{km}$ है। \omterm{prop:g11:seismic-structure:crossover}{अतिक्रमण दूरी} $2 \times 35 \times
\sqrt{14.0/2.0} = 70\sqrt 7 \approx \qty{185}{km}$ है।
\cref{exo:g11:seismic-structure:7} का स्टेशन \qty{200}{km} पर है, \omterm{prop:g11:seismic-structure:crossover}{अतिक्रमण दूरी} के
आगे: \omterm{ex:g11:seismic-structure:sp}{शीर्ष तरंग} पहले पहुँचती है, जैसा पाया गया था।
\end{solution}

\begin{solution}{exo:g11:seismic-structure:9}
$\sin i_c = 6.5/8.0 = 0.8125$, $i_c \approx \ang{54.3}$ और $\cos i_c \approx
0.583$। \omterm{prop:g11:seismic-structure:head-wave}{अंतःखंड काल}: $2 \times 6 \times 0.583/6.5 \approx \qty{1.1}{s}$।
\omterm{prop:g11:seismic-structure:crossover}{अतिक्रमण दूरी}: $12\sqrt{14.5/1.5} \approx \qty{37}{km}$। किसी
महाद्वीप के नीचे \omterm{prop:g11:seismic-structure:crossover}{अतिक्रमण दूरी} \qty{165}{km} के पास होती है, इसलिए कुछ दसियों
किलोमीटर की दूरी पर लगे स्टेशन दोनों रेखाएँ पकड़ लेते हैं; महासागर के नीचे \omterm{ex:g11:seismic-structure:sp}{शीर्ष तरंग} लगभग
\qty{37}{km} के आगे ही पहले पहुँच जाती है, और सीधी तरंग की रेखा केवल
कुछ दसियों किलोमीटर लंबी है: उसे खींचने के लिए स्टेशनों को कुछ किलोमीटर की
दूरी पर होना चाहिए।
\end{solution}

\begin{solution}{exo:g11:seismic-structure:10}
$\sin r = (8.06/13.7)\sin\ang{30} \approx 0.294$, इसलिए $r \approx \ang{17.1}$:
किरण अभिलंब की ओर, यानी पृथ्वी के केंद्र की ओर, मुड़ती है।
इसलिए वह क्रोड को उसके मध्य के पास एक लंबे मार्ग से पार करती है,
और दूसरी ओर बहुत दूर जाकर बाहर निकलती है: क्रोड में प्रवेश करने वाली P किरणें
उन किरणों की अपेक्षा भूकंप से बहुत अधिक दूर निकलती हैं जो उसे बस छूकर निकल जाती हैं।
\end{solution}

\begin{solution}{exo:g11:seismic-structure:11}
$12\,742/13.7 \approx \qty{930}{s}$, लगभग 15.5~मिनट, जबकि तरंग
\qty{1210}{s}, लगभग 20~मिनट, लेती है: वह लगभग 280~सेकंड देर से है। प्रावार की
सबसे अधिक चाल से भी तरंग कहीं पहले पहुँचती, इसलिए
पृथ्वी के मध्य भाग को उसे धीमा करना चाहिए। ओल्डहैम ने निष्कर्ष निकाला कि पृथ्वी का
एक केंद्रीय क्रोड है जिसमें \omterm{def:g11:seismic-structure:body-waves}{P तरंगें} प्रावार की अपेक्षा अधिक धीरे चलती हैं।
\end{solution}

\begin{solution}{exo:g11:seismic-structure:12}
(1) प्रस्ताव के अनुसार $t_i = 2h\cos i_c/v_1$, और $\cos i_c = \sqrt{1 -
v_1^2/v_2^2} = \sqrt{v_2^2 - v_1^2}/v_2$। इसलिए
\[
  t_i = \frac{2h\sqrt{v_2^2 - v_1^2}}{v_1 v_2},
  \qquad\text{अतः}\qquad
  h = \frac{t_i\,v_1 v_2}{2\sqrt{v_2^2 - v_1^2}}.
\]
(2) $h = 6.23 \times 6.2 \times 8.1/\bigl(2\sqrt{65.61 -
38.44}\bigr) \approx 312.9/10.43 \approx \qty{30.0}{km}$।
\end{solution}

\begin{solution}{exo:g11:seismic-structure:13}
(1) $(3480/6371)^3 \approx 0.163$: क्रोड आयतन का लगभग \qty{16}{\percent}
है। (2) माध्य घनत्व आयतन-भारित माध्य है: $5.52 = 0.163
\rho_c + 0.837 \times 4.45$, इसलिए $\rho_c \approx (5.52 - 3.72)/0.163 \approx
\qty{11.0}{g/cm^3}$। (3) यह \qty{9.9}{g/cm^3} से अधिक है, जैसा होना
चाहिए, क्योंकि क्रोड का घनत्व केंद्र की ओर बढ़ता है (\qty{13.1}{g/cm^3}
तक); यह भूपर्पटी के घनत्व, \qty{2.6}{g/cm^3}, का लगभग चार गुना है।
क्रोड सतह की चट्टानों की अपेक्षा कहीं अधिक घने पदार्थ से बना है।
\end{solution}

\begin{solution}{exo:g11:seismic-structure:14}
(1) $105 \times 111.2 \approx \qty{11700}{km}$ और $143 \times 111.2 \approx
\qty{15900}{km}$। (2) $(\cos\ang{105} - \cos\ang{143})/2 \approx (-0.259 +
0.799)/2 \approx 0.27$: सतह का लगभग \qty{27}{\percent}, भूकंप के विपरीत
बिंदु के चारों ओर एक वलयाकार पट्टी। (3) \omterm{def:g11:seismic-structure:shadow-zone}{छाया क्षेत्र} अधिकेंद्र से दूरियों का एक
परिसर है: वह भूकंप के साथ खिसकता है। हर भूकंप की
अपनी वलयाकार पट्टी होती है।
\end{solution}

\begin{solution}{exo:g11:seismic-structure:15}
(1) $h = 9.5 \times 6.3 \times 8.1/\bigl(2\sqrt{65.61 - 39.69}\bigr) \approx
484.8/10.18 \approx \qty{48}{km}$। (2) $\sin i_c = 6.3/8.1 \approx 0.778$,
$\cos i_c \approx 0.629$, और $t_i = 2 \times 30 \times 0.629/6.3 \approx
\qty{6.0}{s}$। (3) श्रेणी के नीचे भूपर्पटी लगभग \qty{48}{km} मोटी है,
जो महाद्वीपीय औसत, लगभग \qty{41}{km}, से अधिक मोटी और पर्वत-पट्टियों के
औसत के निकट है: ऊँची श्रेणियाँ मोटी भूपर्पटी पर खड़ी होती हैं।
\end{solution}

\begin{solution}{pb:g11:seismic-structure:1}
\textbf{1.} Pg बिंदु एक सरल रेखा पर स्थित हैं; Pn बिंदु
दूसरी, कम खड़ी सरल रेखा पर।

\textbf{2.} भूकंप अधिकेंद्र के नीचे बहुत कम गहराई पर है:
शून्य दूरी पर स्थित स्टेशन तरंग को तुरंत दर्ज करता।

\textbf{3.} $(32.3 - 3.3)/(200 - 20) = 29.0/180 \approx \num{0.161}$~\unit{s}
प्रति किलोमीटर।

\textbf{4.} $v_1 = 180/29.0 \approx \qty{6.2}{km/s}$।

\textbf{5.} $3.3 + 280 \times 0.161 \approx \qty{48.4}{s}$।

\textbf{6.} $(43.2 - 18.6)/(300 - 100) = 24.6/200 = \num{0.123}$~\unit{s}
प्रति किलोमीटर, इसलिए $v_2 = 200/24.6 \approx \qty{8.13}{km/s}$।

\textbf{7.} यह \qty{8.0}{km/s} के निकट है: भूपर्पटी के नीचे सबसे ऊपरी
प्रावार की चट्टान है।

\textbf{8.} $\sin i_c = 6.21/8.13 \approx 0.763$, इसलिए $i_c \approx
\ang{49.8}$।

\textbf{9.} \omterm{ex:g11:seismic-structure:sp}{शीर्ष तरंग} कोण $i_c$ पर नीचे जाती है और उसी कोण पर वापस
ऊपर आती है। पहला बिंदु जहाँ वह सतह पर लौट सकती है, उस बिंदु से
$h\tan i_c$ आगे है जहाँ वह \omterm{def:g11:seismic-structure:moho}{मोहो} तक पहुँची थी, और वह बिंदु स्वयं स्रोत से $h\tan i_c$
पर है: $2h\tan i_c$ से निकट कोई स्टेशन उसे प्राप्त नहीं कर सकता।

\textbf{10.} भूकंपमापी अपने तक पहुँचने वाली हर तरंग दर्ज करता है, केवल
पहली नहीं; Pn तरंग Pg तरंग के बाद पहुँचती है, पर अभिलेख पर फिर भी एक अलग
तरंग के रूप में दिखती है।

\textbf{11.} \qty{160}{km} पर Pg (\qty{25.8}{s}) Pn
(\qty{26.0}{s}) से पहले आती है; \qty{180}{km} पर Pn (\qty{28.5}{s}) Pg
(\qty{29.0}{s}) से पहले आती है। बदलाव इन दोनों स्टेशनों के बीच होता है।

\textbf{12.} Pg: $t = 3.3 + 0.1611\,(x - 20) \approx 0.08 + 0.1611\,x$। Pn:
$t = 18.6 + 0.123\,(x - 100) = 6.30 + 0.123\,x$ ($x$ किलोमीटर में, $t$
सेकंड में)।

\textbf{13.} $0.08 + 0.1611\,x = 6.30 + 0.123\,x$ से $x \approx
6.22/0.0381 \approx \qty{163}{km}$ मिलता है।

\textbf{14.} Pn, $6.30 + 0.123 \times 250 \approx \qty{37.1}{s}$ के बाद।

\textbf{15.} $t_i \approx \qty{6.3}{s}$।

\textbf{16.} $\cos i_c \approx 0.646$, इसलिए $h = 6.21 \times 6.3/(2 \times
0.646) \approx \qty{30.3}{km}$।

\textbf{17.} $x_c = 2 \times 30.3 \times \sqrt{14.34/1.92} \approx 60.6 \times
2.73 \approx \qty{165}{km}$, प्रश्न~13 के \qty{163}{km} से
\qty{2}{km} के भीतर।

\textbf{18.} $\tan i_c \approx 1.18$: $2 \times 30.3 \times 1.18 \approx
\qty{72}{km}$। Pn दर्ज करने वाला सबसे निकट का स्टेशन \qty{80}{km} पर है, और
\qty{60}{km} वाला उसे दर्ज नहीं करता: दोनों मेल खाते हैं।

\textbf{19.} समय \qty{0.1}{s} तक पढ़े जाते हैं, और केवल $t_i$ पर
\qty{0.1}{s} की त्रुटि $h$ को लगभग \qty{0.5}{km} बदल देती है; ढालों से
ली गई चालें भी अनिश्चित हैं, और भूकंप ठीक सतह पर नहीं
था। भूपर्पटी भी पूरी तरह एकसमान नहीं है। एक-दो किलोमीटर की परिशुद्धता ही
अधिकतम है जिसकी आँकड़े अनुमति देते हैं।

\textbf{20.} \textbf{इस क्षेत्र के नीचे \omterm{def:g11:seismic-structure:moho}{मोहो} लगभग \qty{30}{km} पर है}:
यहाँ की भूपर्पटी किसी महाद्वीप के लिए पतली ओर है, जिसकी भूपर्पटी औसतन लगभग
\qty{41}{km} होती है; \qty{30}{km} उन क्षेत्रों की विशेषता है जहाँ भूपर्पटी
खिंची है।
\end{solution}
